\subsubsection{Segment Tree Lazy Assign + Sum}
\begin{lstlisting}[style=mystyle, language=C++]
// TC: O(log n) por operacion, O(n) en memoria
// usa: asignar un valor a un rango y consultar la suma en un rango
struct SegmentTree {
	struct Node {
		int val = 0;
		int lazy = LLONG_MAX; // valor neutro
	};

	int n;
	vector<Node> st;

	SegmentTree(int size) {
		n = size;
		while(__builtin_popcount(n) != 1) n++;
		st.resize(2*n);
	}

	void push(int node, int l, int r) {
		if(st[node].lazy == LLONG_MAX) return;
		st[node].val = (r-l+1) * st[node].lazy;
		if(l != r) {
			st[node*2].lazy = st[node].lazy;
			st[node*2+1].lazy = st[node].lazy;
		}
		st[node].lazy = LLONG_MAX;
	}

	void assign_range(int node, int l, int r, int ql, int qr, int val) {
		push(node, l, r);
		if(ql <= l && r <= qr) {
			st[node].lazy = val;
			push(node, l, r);
			return;
		}
		if(r < ql || qr < l) return;
		int mid = l + (r-l)/2;
		assign_range(node*2, l, mid, ql, qr, val);
		assign_range(node*2+1, mid+1, r, ql, qr, val);
		st[node].val = st[node*2].val + st[node*2+1].val;
	}

	int query(int node, int l, int r, int ql, int qr) {
		push(node, l, r);
		if(ql <= l && r <= qr) return st[node].val;
		if(r < ql || qr < l) return 0;
		int mid = l + (r-l)/2;
		return query(node*2, l, mid, ql, qr)
		+ query(node*2+1, mid+1, r, ql, qr);
	}

	void assign(int l, int r, int val) {
		assign_range(1, 0, n-1, l, r, val);
	}
	int sum(int l, int r) {
		return query(1, 0, n-1, l, r);
	}
};
int main(){
	SegmentTree st(5);
	st.assign(0, 4, 1);          // [1,1,1,1,1]
	st.assign(1, 3, 2);          // [1,2,2,2,1]
	cout << st.sum(0, 4) << endl; // 8
	cout << st.sum(1, 3) << endl; // 6
	st.assign(0, 4, 5);          // [5,5,5,5,5]
	cout << st.sum(0, 4) << endl; // 25
	return 0;
}
\end{lstlisting}
